\hypertarget{methodology-project-development}{%
\section{3. Methodology / Project
Development}\label{methodology-project-development}}

\begin{quote}
The bulk of the thesis. Three coordinated work-streams: \textbf{3.A} the
network, \textbf{3.B} the endpoints, \textbf{3.C} the knowledge
artefact. Maps to LaTeX section \emph{Methodology / project
development}.
\end{quote}

\begin{center}\rule{0.5\linewidth}{0.5pt}\end{center}

\hypertarget{a-network-hardware-deployment}{%
\subsection{3.A Network hardware
deployment}\label{a-network-hardware-deployment}}

This part of the chapter describes the hardware-network work-stream: how
the physical and link layers of a small community network are designed,
specified, and brought up in a low-resource site. The narrative is
deliberately ordered as a deployment would be ordered in the field, but
it is grouped under a four-layer model --- \emph{field}, \emph{site},
\emph{endpoint touchpoint}, \emph{power and enclosure} --- borrowed from
the functional layering proposed in the companion thesis {[}Motje,
2026{]}. Adopting the same layer vocabulary on the hardware side and on
the software/services side keeps the two theses, and the handbook that
backs them, internally coherent: every recipe in \texttt{docs/3-Guide/}
of the handbook can be located in exactly one layer, and every artefact
named below has a one-to-one counterpart in the handbook.

The material is also organised around three validation lenses that recur
in §4: \textbf{coverage} (does the layer reach where it needs to
reach?), \textbf{sufficiency} (does it carry the expected load with
margin?), and \textbf{adaptation} (does the recipe survive being
executed by a different team in a different site?). Where a design
choice trades one lens against another, the trade-off is named
explicitly.

\hypertarget{a.1-the-four-layer-model}{%
\subsubsection{3.A.1 The four-layer
model}\label{a.1-the-four-layer-model}}

The hardware deployment is decomposed into four functional layers. Each
layer owns a small set of decisions, a recommended bill of materials,
and a matching family of recipes in the handbook.

\begin{longtable}[]{@{}
  >{\raggedright\arraybackslash}p{(\columnwidth - 6\tabcolsep) * \real{0.2500}}
  >{\raggedright\arraybackslash}p{(\columnwidth - 6\tabcolsep) * \real{0.2500}}
  >{\raggedright\arraybackslash}p{(\columnwidth - 6\tabcolsep) * \real{0.2500}}
  >{\raggedright\arraybackslash}p{(\columnwidth - 6\tabcolsep) * \real{0.2500}}@{}}
\toprule\noalign{}
\begin{minipage}[b]{\linewidth}\raggedright
Layer
\end{minipage} & \begin{minipage}[b]{\linewidth}\raggedright
Owns
\end{minipage} & \begin{minipage}[b]{\linewidth}\raggedright
Handbook section
\end{minipage} & \begin{minipage}[b]{\linewidth}\raggedright
Examples of artefacts
\end{minipage} \\
\midrule\noalign{}
\endhead
\bottomrule\noalign{}
\endlastfoot
\textbf{Field} & Inter-building radio links, long-distance bridges,
outdoor antennas & \texttt{3-Guide/Antennas/} & Point-to-point Ubiquiti
links, line-of-sight surveys \\
\textbf{Site} & Indoor coverage, mesh backbone, IP plan, switching,
boundary L3/L4 services & \texttt{3-Guide/Network-Planning/},
\texttt{IP-Addressing/}, \texttt{Wireless-Mesh/} & OpenWrt routers,
802.11s mesh, DHCP/DNS on the gateway \\
\textbf{Endpoint touchpoint} & The network-side of laptop/desktop
provisioning & \texttt{3-Guide/Laptop-Deployment/} (network parts only)
& Isolated PXE subnet, deployment switch, on-site DHCP/TFTP/NFS \\
\textbf{Power \& enclosure} & Mains, UPS, surge protection, mounting,
cabling, labelling & \texttt{3-Guide/Power-and-UPS/} & UPS sizing, cable
management, panel labelling \\
\end{longtable}

The endpoint-touchpoint layer is intentionally split from the
\emph{endpoint} work-stream described in §3.B. Section §3.B owns the
laptop as a managed asset (image, partitions, user account); §3.A only
owns the network plumbing that is required during a mass deployment. The
split mirrors the SW/HW split between this thesis and the companion
thesis: a deployment is a single event, but the underlying
responsibilities can be assigned cleanly.

The remainder of §3.A walks through these four layers in the order in
which they are typically commissioned: a \emph{site assessment} that
constrains all later choices (§3.A.2), the \textbf{site layer} itself
(§3.A.3 to §3.A.5), the \textbf{field layer} (§3.A.6), the
\textbf{endpoint touchpoint} (§3.A.7), the \textbf{power and enclosure
layer} (§3.A.8), and the \textbf{boundary services} that are kept in
scope (§3.A.9). The part closes with the cross-cutting practices that
experience showed are not optional (§3.A.10) and the explicit
lessons-to-handbook flow that motivated the recipes (§3.A.11).

\hypertarget{a.2-site-and-internet-assessment}{%
\subsubsection{3.A.2 Site and internet
assessment}\label{a.2-site-and-internet-assessment}}

No layer can be designed without first characterising the site. The
method codified in \texttt{3-Guide/Network-Planning/} has three phases,
each with a corresponding deliverable:

\begin{enumerate}
\def\labelenumi{\arabic{enumi}.}
\tightlist
\item
  \textbf{Internet assessment.} What ISP options exist, what speed and
  latency the existing uplink actually delivers, and whether a different
  technology (fibre, cable, ADSL, 4G/5G, LEO satellite) is reachable at
  reasonable cost. The deliverable is a comparison table scored against
  the eight criteria in \texttt{1-Internet-Assessment.md}
  (download/upload, data cap, reliability, latency, local support,
  contract length, budget fit, expected user load). The Namibia
  deployment inherited a government-supplied ADSL link and skipped the
  change-of-provider branch; the methodology still required documenting
  the alternatives, which is what allows a future team to revisit the
  choice.
\item
  \textbf{Site assessment.} A walkthrough with a Wi-Fi analyser
  (LinSSID/WiFiman) producing a coverage heatmap with three bands (≥ −50
  dBm, −70 to −80 dBm, ≤ −80 dBm), a baseline speed test repeated at
  different times of day, and an annotated site map (overhead imagery
  plus on-site sketch) marking entry point, target rooms, obstacles,
  distances, and power outlets. This map is the single most important
  artefact produced before any equipment is purchased.
\item
  \textbf{Expansion planning.} A decision tree (Ethernet + APs vs mesh
  vs point-to-point antennas) driven by \emph{building topology, cabling
  feasibility, and inter-building distance}, not by personal preference.
\end{enumerate}

The methodological commitment is that the site survey produces a written
record, not a recollection. Every later decision in §3.A.3 to §3.A.8
cites this record. This is the first concrete answer to the
\emph{knowledge volatility} problem named in §1.2: the site assessment
of one team becomes the starting context of the next.

\hypertarget{a.3-hardware-selection-at-the-site-layer}{%
\subsubsection{3.A.3 Hardware selection at the site
layer}\label{a.3-hardware-selection-at-the-site-layer}}

The bill of materials for the site layer is constrained by four
criteria, applied in order:

\begin{enumerate}
\def\labelenumi{\arabic{enumi}.}
\tightlist
\item
  \textbf{OpenWrt support} with mesh-capable Wi-Fi packages. The router
  must be on the OpenWrt Table of Hardware with a current snapshot or
  release, and it must have enough flash and RAM to host
  \texttt{wpad-mesh-wolfssl} alongside the default LuCI stack (in
  practice ≥ 16 MB flash, ≥ 128 MB RAM, dual-band radios).
\item
  \textbf{Affordability and local availability.} Equipment that can be
  replaced in country, or shipped without customs friction, beats
  marginally better equipment that cannot.
\item
  \textbf{Power flexibility.} 12 V / USB-C devices integrate with a
  small UPS without a per-router AC brick.
\item
  \textbf{Footprint.} A consumer plastic enclosure that mounts on a wall
  with two screws is preferred over a 19'' rack device for
  primary-school environments.
\end{enumerate}

The default pick that satisfies all four criteria at the time of writing
is the \textbf{Cudy WR3000E} (Wi-Fi 6, dual-band, OpenWrt-supported,
\textasciitilde45 €/unit, 12 V input). For deployments that need a more
capable gateway with two Ethernet interfaces, USB 3.0, and headroom for
VPN and on-router services, the \textbf{NanoPi R-series} is the chosen
alternative; it is the same family that the companion thesis uses for
service hosting {[}Motje, 2026{]}, which keeps the on-site hardware
vocabulary uniform across the two work-streams.

A pitfall worth naming: a router that \emph{technically} supports
OpenWrt at release N may be impractical if the snapshot images regress
on the model's specific Wi-Fi driver. Two days were lost in early
bring-up to a snapshot build that broke 5 GHz on the WR3000E; the recipe
in \texttt{Wireless-Mesh/1-Static-IP-Mesh/index.md} therefore pins the
tested versions of OpenWrt and \texttt{wpad-mesh-wolfssl} and cautions
the reader that ``newer is not necessarily better''.

\hypertarget{a.4-ip-addressing-plan}{%
\subsubsection{3.A.4 IP addressing plan}\label{a.4-ip-addressing-plan}}

The IP plan is the most under-valued artefact of a small network, and
the one whose absence costs the most when something breaks. The handbook
recipe \texttt{3-Guide/IP-Addressing/index.md} codifies four
methodological commitments:

\begin{itemize}
\tightlist
\item
  \textbf{Use a non-default \texttt{/24}} (the deployments documented
  here use \texttt{192.168.70.0/24}). Default ranges
  (\texttt{192.168.1.0/24}, \texttt{192.168.0.0/24}) collide the moment
  a second consumer router is plugged in for tests, or a contractor's
  device with the same default range is brought to site. RFC 1918 leaves
  an enormous private space; using it removes a recurring source of
  accidental conflicts.
\item
  \textbf{Reserve the low part of the range for infrastructure}
  (\texttt{.1} for the gateway, \texttt{.2}--\texttt{.20} for routers
  and APs, \texttt{.21}--\texttt{.99} for fixed servers and printers)
  and the high part for DHCP-assigned client devices
  (\texttt{.100}--\texttt{.200}). The split is arbitrary but stable, and
  a single glance at an IP tells the operator what kind of device they
  are dealing with.
\item
  \textbf{Document the plan in a versioned spreadsheet} with one row per
  device: hostname, role, location, MAC, IP, allocation method (static /
  DHCP reservation / DHCP pool). The spreadsheet lives in the project
  repository and is updated \emph{at the moment of provisioning}, not
  afterwards.
\item
  \textbf{Prefer DHCP reservations over hand-coded static IPs} wherever
  the device supports DHCP-client mode. This is the substance of the
  second iteration of the mesh setup (§3.A.5): static IPs survive only
  as long as the team that wrote them is around, while a DHCP table on
  the gateway is self-documenting and survives turnover.
\end{itemize}

The IP-conflict episode of Day 8 in Gochas, in which a newly added
router silently misrouted traffic until a missing static-route entry was
found, is the empirical justification for these commitments. It cost
roughly an hour of debugging for what was a five-minute fix once the
problem was named --- the handbook recipe explicitly calls this out
under ``Common Mistakes'' so that the next operator skips the hour.

\hypertarget{a.5-wireless-mesh-design-at-the-site-layer}{%
\subsubsection{3.A.5 Wireless mesh design at the site
layer}\label{a.5-wireless-mesh-design-at-the-site-layer}}

Indoor coverage of a multi-room or multi-building site without trenching
cable is met with an IEEE 802.11s mesh on OpenWrt. The handbook codifies
the design as \textbf{two iterations} that are meant to be done in
sequence, not chosen between (\texttt{3-Guide/Wireless-Mesh/index.md}):

\begin{itemize}
\tightlist
\item
  \textbf{Iteration 1 --- static-IP mesh.} Each satellite router gets a
  unique static LAN IP, its DHCP server is disabled, the default
  \texttt{wpad-basic-*} package is replaced with
  \texttt{wpad-mesh-wolfssl}, and a 5 GHz 802.11s backhaul is brought up
  alongside a shared 2.4 GHz access point. This is the simplest
  configuration that confirms the mesh links form, that satellites
  bridge correctly, and that clients roam between APs sharing the same
  ESSID, encryption, and key. Iteration 1 is the recommended starting
  point for \emph{every} deployment because it isolates the
  link-formation problem from the IP-management problem.
\item
  \textbf{Iteration 2 --- DHCP-based mesh.} Once the mesh is stable,
  each satellite is converted to a DHCP client of the main router, with
  WAN and the firewall disabled, and its address pinned by a static
  lease on the gateway. This centralises IP management on one device,
  makes onboarding new satellites a matter of plugging them in, and
  removes a whole class of long-term operational problems caused by
  hand-edited IP tables drifting out of date.
\end{itemize}

The methodological choice of \emph{layering} the iterations rather than
presenting them as alternatives is itself a deliberate didactic device:
new operators learn the system by building it incrementally, and they
acquire diagnostic intuition (what fails when DHCP is misconfigured vs
when the mesh ID is mistyped) that they would not acquire by
copy-pasting the final state.

The radio design splits the two bands by role:

\begin{itemize}
\tightlist
\item
  \textbf{5 GHz} carries the mesh backhaul on a fixed channel (typically
  44), 20 MHz or 40 MHz wide, in 802.11s mode with WPA3-SAE encryption.
  Narrower channels penetrate walls better at the cost of throughput; in
  small primary-school deployments the throughput ceiling is set by the
  uplink, not the backhaul, so 20 MHz is the safe default.
\item
  \textbf{2.4 GHz} carries the client access point on a non-overlapping
  channel (1, 6, or 11) with WPA2-PSK by default, escalating to
  WPA2/WPA3 mixed mode only after the client device population is known.
  The Gochas deployment hit a WPA3-only incompatibility on a single
  laptop on Day 6 and was reconfigured to mixed mode within minutes; the
  handbook now recommends mixed mode by default for any site with
  unknown client hardware.
\end{itemize}

\hypertarget{a.6-field-layer-point-to-point-links}{%
\subsubsection{3.A.6 Field-layer point-to-point
links}\label{a.6-field-layer-point-to-point-links}}

When two buildings are more than \textasciitilde50 m apart or separated
by structures that block 2.4/5 GHz, the mesh stops being adequate and a
dedicated point-to-point (PtP) link at the field layer takes over. The
handbook section \texttt{3-Guide/Antennas/} is currently a stub and is
one of the deliberate WIP markers acknowledged in §1.4. The
methodological intent, already fixed, is:

\begin{itemize}
\tightlist
\item
  Use Ubiquiti airMAX or LiteBeam-class equipment (or equivalent) chosen
  by Ubiquiti Design Center / LinkPlanner against a real terrain profile
  rather than a rule of thumb.
\item
  Conduct an explicit line-of-sight visibility test before mounting,
  including Fresnel-zone clearance.
\item
  Mount, align, and link-test in two passes --- coarse alignment from
  the signal-strength meter, fine alignment from the throughput test.
\item
  Document the link as a pair of devices with their own subnet entries
  in the IP plan of §3.A.4.
\end{itemize}

The Gochas deployment did not exercise this layer (the school's
buildings are within mesh range), so the recipe is documented from the
network planning context and from the literature rather than from a
fresh field event. This honest framing is part of the validation
discussion of §4.

\hypertarget{a.7-endpoint-touchpoint-the-deployment-subnet}{%
\subsubsection{3.A.7 Endpoint touchpoint --- the deployment
subnet}\label{a.7-endpoint-touchpoint-the-deployment-subnet}}

A laptop mass-deployment (§3.B) requires its own short-lived network: an
isolated Ethernet segment carrying DHCP, TFTP, and NFS for PXE boot,
plus the Clonezilla rootfs and the disk image. From the §3.A
perspective, this network has three hardware requirements:

\begin{enumerate}
\def\labelenumi{\arabic{enumi}.}
\tightlist
\item
  A \textbf{gigabit unmanaged switch} with one port per target machine
  plus one port for the PXE server. Nine ThinkPads in Gochas were served
  by a single 16-port switch.
\item
  \textbf{Strict isolation from the production LAN} during the imaging
  window. The deployment subnet runs its own DHCP server (the PXE
  server's \texttt{isc-dhcp-server}) which would conflict with the
  gateway's DHCP if both were reachable on the same broadcast domain.
\item
  \textbf{A planned IP range that does not overlap the production plan}
  (\texttt{10.10.10.0/24} was used in Gochas against the production
  \texttt{192.168.70.0/24}), so a misconnected cable does not silently
  bridge the two networks.
\end{enumerate}

The corresponding services (DHCP scope, TFTP root, NFS export) belong to
§3.B; their hardware substrate belongs here.

\hypertarget{a.8-power-and-enclosure}{%
\subsubsection{3.A.8 Power and
enclosure}\label{a.8-power-and-enclosure}}

Power is the failure mode that no software can mitigate. The site
assessment of §3.A.2 records a load inventory for each candidate
deployment site:

\begin{itemize}
\tightlist
\item
  Mesh router: \textasciitilde5 W steady, \textasciitilde10 W peak (per
  node).
\item
  PoE switch (if used): \textasciitilde10--20 W plus per-port PoE
  budget.
\item
  Mini-PC / NanoPi gateway: \textasciitilde5--15 W steady.
\item
  Optional UPS standby loss: \textasciitilde5 W.
\end{itemize}

At a small site (one gateway, one switch, four to seven mesh nodes), the
total steady draw is in the 50--100 W range, which puts a 600--800 VA
consumer UPS in the right size class for \textasciitilde30 min of
bridging during short outages and, more importantly, for graceful
shutdown on longer outages. The recurring power cuts in Gochas
(sometimes lasting days, with the town's communication tower turning off
at 20:00 by design) confirmed that \emph{riding through} outages is not
a realistic goal at this site class; the realistic goal is
\emph{graceful shutdown and clean restart}, which the chosen sizing
supports. The companion thesis {[}Motje, 2026{]} develops the
graceful-shutdown integration on the services side via NUT; §3.A only
guarantees that the hardware can survive a sudden cut without data loss
on the network devices themselves (OpenWrt's flash layout makes this
trivially true on the router side; the gateway PC is the device that
needs the UPS).

Enclosure and mounting are mundane and decisive. The convention adopted
in the handbook recipe (currently a stub at
\texttt{3-Guide/Power-and-UPS/}) is: mount above head height to clear
furniture and people; choose locations with a power outlet within
\textasciitilde2 m to avoid extension-cord chains; label every cable on
both ends with a printed tag bearing the device hostname; photograph the
panel after every change.

\hypertarget{a.9-boundary-services-kept-in-scope}{%
\subsubsection{3.A.9 Boundary services kept in
scope}\label{a.9-boundary-services-kept-in-scope}}

The hardware work-stream owns the \emph{minimum} set of L3/L4 services
without which the network is not usable. Three services qualify:

\begin{itemize}
\tightlist
\item
  \textbf{DHCP on the gateway.} A single \texttt{dnsmasq} instance (the
  OpenWrt default) serving the production scope, with static leases for
  every infrastructure device named in §3.A.4. The SW/HW boundary is
  drawn here: configuration of \texttt{dnsmasq} on the gateway is in
  scope; richer DHCP servers, DHCP relays across VLANs, and IPAM tooling
  are in {[}Motje, 2026{]}.
\item
  \textbf{DNS forwarding on the gateway.} Again \texttt{dnsmasq} in its
  forwarder role, with the gateway's own DNS servers (typically
  \texttt{1.1.1.1} and \texttt{9.9.9.9}) as upstream. Local-name
  resolution for the infrastructure devices is configured by the same
  static-lease entries.
\item
  \textbf{A monitoring overview on the gateway.} Only the \emph{minimum}
  needed to answer ``is the mesh up, are the satellites reachable, is
  the uplink alive'': LuCI's built-in \texttt{Status\ →\ Wireless} and
  \texttt{Status\ →\ Routes}, plus ping-based checks scriptable from the
  gateway. Full Zabbix agentful monitoring is software work and is owned
  by {[}Motje, 2026{]}; the hardware-side commitment is to expose the
  SNMP and ping endpoints that the monitoring stack consumes.
\end{itemize}

This boundary is drawn explicitly so that no aspect of the deployment
falls between the two theses uncovered.

\hypertarget{a.10-cross-cutting-operational-practices}{%
\subsubsection{3.A.10 Cross-cutting operational
practices}\label{a.10-cross-cutting-operational-practices}}

A handful of practices are not specific to any one layer but apply
across all of them. They are codified in the handbook's general rules
and repeated in the site-deployment recipes because their cost of
omission is high:

\begin{itemize}
\tightlist
\item
  \textbf{Label every cable on both ends.} A two-letter site code plus
  the destination hostname is enough.
\item
  \textbf{Photograph every panel} before and after a change. The photo
  is attached to the project log entry.
\item
  \textbf{Keep a one-page outage runbook} at the site (laminated A4),
  containing: gateway IP and admin credentials reset procedure, the IP
  plan, the wiring diagram, and the contact of the maintainer of record.
\item
  \textbf{Carry a spares kit} sized to one of each critical device (one
  router, one switch, one PoE injector if applicable, two pre-made
  Ethernet cables of common lengths). The kit lives at the site, not in
  Barcelona.
\item
  \textbf{Pin firmware and package versions} in the recipe and verify
  them during bring-up. The ``newer is not necessarily better'' rule of
  §3.A.3 applies here too.
\end{itemize}

These practices are mundane on purpose. Their value is exactly that they
do not require the operator to remember anything in the moment.

\hypertarget{a.11-lessons-learned-and-inputs-to-the-handbook}{%
\subsubsection{3.A.11 Lessons learned and inputs to the
handbook}\label{a.11-lessons-learned-and-inputs-to-the-handbook}}

The hardware-network deployments described in §4 produced a set of
lessons that landed directly as recipes, admonitions, or
cross-references in the handbook. Naming them here makes the §3 →
handbook → §4 loop explicit.

\textbf{The IP plan is a deliverable, not a side-effect.} The Day-8
conflict episode in Gochas is the reason
\texttt{3-Guide/IP-Addressing/index.md} opens with a ``why a non-default
range matters'' section and why every router recipe asks the reader to
update the spreadsheet \emph{at the moment of provisioning}.

\textbf{Mesh setup is two iterations, not one.} The decision to split
the mesh recipe into a static-IP first pass and a DHCP-based second pass
came directly from the observation that operators who tried to go to the
DHCP-based design first lost time on link-formation failures that were
not in fact caused by DHCP. The two recipes
(\texttt{Wireless-Mesh/1-Static-IP-Mesh/},
\texttt{Wireless-Mesh/2-DHCP-Mesh/}) are the codification of that
lesson.

\textbf{Pin the OpenWrt and \texttt{wpad-mesh} versions.} The
\texttt{Used\ Versions} table at the top of the static-IP-mesh recipe
exists because two days of early bring-up were lost to a snapshot
regression. Future operators are warned that they may need to roll back
if a fresh snapshot misbehaves, and they are pointed to the issue
tracker.

\textbf{Mixed-mode WPA2/WPA3 by default until clients are known.} The
Day-6 WPA3-only incompatibility is named in the recipe under ``Wireless
Security'' so the next deployment does not repeat it.

\textbf{Power planning is for graceful shutdown, not ride-through.} The
recurring outages in Gochas reframed the UPS goal and that reframing is
now embedded in the (still stub) \texttt{3-Guide/Power-and-UPS/}
section. The handover to {[}Motje, 2026{]} for NUT integration is
explicit.

\textbf{The site survey is a versioned artefact.} The methodological
commitment that ``every later decision cites this record'' is inherited
from §3.A.2 and is repeated as a rule in
\texttt{Network-Planning/index.md}. A site survey kept only in someone's
head does not count.

These lessons are revisited in §4 against the three validation lenses
(coverage, sufficiency, adaptation) and in §7 as part of the
per-objective check.

\begin{center}\rule{0.5\linewidth}{0.5pt}\end{center}

\hypertarget{b-endpoint-reconditioning-and-mass-deployment}{%
\subsection{3.B Endpoint reconditioning and mass
deployment}\label{b-endpoint-reconditioning-and-mass-deployment}}

This part of the chapter describes the second hardware work-stream: how
a batch of refurbished laptops is taken from incoming-equipment status
to classroom-ready endpoints with a single, reproducible image. The
narrative is again ordered as a deployment is ordered in the field, but
the underlying methodological commitments are the same as those of §3.A:
written artefacts over recollection, recipes that survive a change of
operator, and explicit lessons looped back into the handbook.

§3.B owns the laptop as a managed asset (the image, the partitions, the
user account, the BIOS posture). The network plumbing that the
deployment briefly needs (the isolated PXE subnet, the deployment
switch, the DHCP/TFTP/NFS hardware footprint) is owned by §3.A.7.
Anything beyond first boot --- fleet management, identity, application
updates, remote support --- belongs to the companion thesis {[}Motje,
2026{]}.

\hypertarget{b.1-the-refurbished-hardware-case}{%
\subsubsection{3.B.1 The refurbished-hardware
case}\label{b.1-the-refurbished-hardware-case}}

The starting point is the observation that the bottleneck for digital
inclusion in low-resource sites is rarely \emph{new} hardware. Corporate
fleet refreshes, NGOs such as \textbf{Labdoo}, and local sponsors put
three- to seven-year-old business laptops back into circulation by the
container-load; the limiting factor on their reuse is not the silicon
but the \emph{labour cost of imaging them one by one} and the \emph{cost
of a support contract for a fleet of mismatched configurations}. The
methodological response of this work-stream is therefore not ``find
better hardware'' but ``drive the per-machine provisioning cost as close
to zero as possible while keeping the resulting fleet uniform enough to
support remotely''.

Two sourcing channels were exercised in the deployments documented in
§4. Labdoo provided nine Lenovo ThinkPads (T460 and X260, Intel
i5-6200U, 8 GB DDR4, mixed 238 GB SSD / 466 GB HDD) wiped and ready to
receive a new OS. NexTReT contributed three additional laptops and two
mini-PC servers from a Spanish fleet refresh. Both channels deliver
hardware in the \emph{same generation class} but with
\emph{non-identical disk geometries} --- which, as §3.B.5 will show, is
the single decision driver of the entire imaging workflow.

The case for refurbished hardware is also environmental. A laptop
manufactured five years ago and reused for another five years has a
markedly lower lifecycle carbon footprint per useful year than a
newly-manufactured equivalent; this argument is developed in §6.

\hypertarget{b.2-intake-triage-and-inventory}{%
\subsubsection{3.B.2 Intake, triage and
inventory}\label{b.2-intake-triage-and-inventory}}

Before any laptop receives the golden-master image, it goes through a
short triage step that produces an inventory record. The fields captured
per machine are deliberately minimal so the workflow does not become its
own bottleneck:

\begin{longtable}[]{@{}
  >{\raggedright\arraybackslash}p{(\columnwidth - 4\tabcolsep) * \real{0.3333}}
  >{\raggedright\arraybackslash}p{(\columnwidth - 4\tabcolsep) * \real{0.3333}}
  >{\raggedright\arraybackslash}p{(\columnwidth - 4\tabcolsep) * \real{0.3333}}@{}}
\toprule\noalign{}
\begin{minipage}[b]{\linewidth}\raggedright
Field
\end{minipage} & \begin{minipage}[b]{\linewidth}\raggedright
Source
\end{minipage} & \begin{minipage}[b]{\linewidth}\raggedright
Why
\end{minipage} \\
\midrule\noalign{}
\endhead
\bottomrule\noalign{}
\endlastfoot
Manufacturer / model / serial & Sticker or
\texttt{dmidecode\ -s\ system-serial-number} & Traceability, warranty \\
CPU / RAM & \texttt{lscpu}, \texttt{free\ -h} & Software-baseline
check \\
Storage device + size & \texttt{lsblk\ -d\ -o\ NAME,SIZE,ROTA} &
\textbf{Drives the imaging plan} \\
Battery health &
\texttt{upower\ -i\ \$(upower\ -e\ \textbackslash{}\textbar{}\ grep\ BAT)}
& Field viability \\
BIOS posture & Manual: Secure Boot off, USB boot on, network boot on &
Pre-deployment requirement \\
Notes & Free text & Visible defects, dead keys, screen marks \\
\end{longtable}

The inventory is held in the same project repository as the IP plan of
§3.A.4, in a comma-separated file with one row per machine. The
deployment script consumes this file when matching disks to image
variants (§3.B.5). For larger or longer-running operations, a tool such
as \textbf{DeviceHub} can replace the spreadsheet without changing the
methodology --- what matters is that the inventory exists and is
versioned, not which tool produces it.

The triage step also enforces a one-time \textbf{BIOS posture} that
every target machine must reach before it joins the deployment queue:
Secure Boot disabled, network boot enabled in the boot order, USB boot
enabled as a fallback. This posture is why §3.B.7 succeeds at scale;
skipping it is the single most expensive mistake an inexperienced
operator can make (see the lessons of §3.B.9).

\hypertarget{b.3-the-golden-master-image-aucoop-linux-mint}{%
\subsubsection{3.B.3 The golden-master image --- AUCOOP Linux
Mint}\label{b.3-the-golden-master-image-aucoop-linux-mint}}

The reference operating system installed on every endpoint is
\textbf{Linux Mint 22.3 Cinnamon}, customised for community use. The
reasoning is documented in
\texttt{3-Guide/Laptop-Deployment/AUCOOP-image.md} and reduces to three
observations.

First, the alternative operating systems each fail one constraint that
matters for this user population. Windows is licensed and runs poorly on
the target hardware class; Ubuntu's GNOME and Snap-centric desktop is
unfamiliar to users coming from Windows and increases first-use
friction; rolling-release distributions impose a maintenance burden that
the receiving institution cannot absorb. Linux Mint Cinnamon presents a
Windows-style desktop (start menu in the bottom-left, taskbar with
launchers, system tray on the bottom-right) on top of an LTS Ubuntu
base, which keeps user training cost and security-update cadence both
manageable.

Second, the \emph{customisation} is itself a methodological deliverable,
not a matter of taste. The AUCOOP image ships with \textbf{OnlyOffice}
as the office suite (chosen for its high fidelity to the Microsoft
Office file formats teachers already produce), with familiar launcher
icons named after the closest Microsoft equivalent (Word, Excel,
PowerPoint), and with the default unused applications removed. The user
account (\texttt{aucoop}, with a documented password) is generic on
purpose: the endpoint is delivered as a \emph{station} the school can
assign to a child or a teacher, not as a personal device tied to the
contributor who imaged it.

Third, the master is captured with a strict pre-capture cleanup
checklist: \texttt{apt\ clean}, \texttt{apt\ autoremove}, removal of
thumbnail caches, truncation of \texttt{journalctl} logs, removal of the
network-manager connection history, and a final fill-with-zeros of the
free space so gzip compression of the captured image is dense. Skipping
the cleanup inflates the image by a factor of two to three with no
functional gain.

\hypertarget{b.4-image-capture-with-clonezilla}{%
\subsubsection{3.B.4 Image capture with
Clonezilla}\label{b.4-image-capture-with-clonezilla}}

The captured artefact is a \textbf{Clonezilla image} of the master disk,
produced from a Clonezilla Live USB booted on the master machine. The
operative choices are:

\begin{itemize}
\tightlist
\item
  \texttt{device-image} mode: the source is a block device, the
  destination is an image directory.
\item
  \texttt{local\_dev} repository: an external USB drive or HDD
  physically separate from the source disk (writing the image to the
  same disk being read is unsupported and would corrupt the result).
\item
  \texttt{savedisk}: the entire disk is captured, not just the root
  partition, so the GPT, the EFI system partition, and the root
  partition are all preserved as a self-contained set.
\item
  gzip compression at the default level: a typical Linux Mint master
  with \textasciitilde12 GB of used data on a 466 GB disk produces a
  \textasciitilde4 GB compressed image directory.
\end{itemize}

The image directory layout is meaningful. It contains one
\texttt{*-ptcl-img.gz} file per partition (e.g.
\texttt{sda1.vfat-ptcl-img.gz}, \texttt{sda2.ext4-ptcl-img.gz}),
partition-table dumps in \texttt{parted} and \texttt{sgdisk} formats
(\texttt{sda-pt.parted}, \texttt{sda-gpt.sgdisk}), a \texttt{parts}
listing, and a \texttt{disk} descriptor. The recovery side (§3.B.7)
reconstructs the geometry from the dumps before restoring the partition
data. Treating the image as a \emph{named, versioned artefact} with a
date-bearing directory name (\texttt{aucoop-mint22.3-2026-03}) is what
makes the per-deployment provenance trail of §4 possible.

\hypertarget{b.5-the-partition-resize-problem-and-the-fix}{%
\subsubsection{3.B.5 The partition-resize problem (and the
fix)}\label{b.5-the-partition-resize-problem-and-the-fix}}

The technical core of this work-stream is a problem that does not appear
when all target disks are the same size and is unavoidable when they are
not. It is described here in full because (i) it has cost more
field-debugging time than any other issue in the deployments documented
in §4, and (ii) it is the kind of failure that mainstream tutorials skip
and on-site teams therefore re-discover the hard way.

\textbf{The symptom.} A Clonezilla image captured from a 466 GB HDD and
restored to a 238 GB SSD fails roughly 77 \% of the way through with:

\begin{verbatim}
target seek ERROR: Invalid argument
\end{verbatim}

\textbf{The cause.} Clonezilla preserves the source partition layout in
the image. When restoring, \texttt{partclone} writes data blocks at the
same offsets at which they appear in the source filesystem. The ext4
filesystem scatters its metadata --- block-group descriptors, inode
tables, block bitmaps, inode bitmaps --- across the \emph{entire}
partition, not only the populated region. A 466 GB ext4 partition
therefore has metadata blocks at offsets up to \textasciitilde466 GB
even when the filesystem holds only 12 GB of data. When the target
partition is smaller than the source \emph{partition} (not the source
\emph{data}), the seek to the high-offset metadata block lands beyond
the device boundary and \texttt{partclone} aborts.

\textbf{Why the obvious flags are not enough.} Two flags from the
\texttt{ocs-sr} manual look as if they should solve this: \texttt{-k1}
(proportionally resize target partitions) and \texttt{-icds} (skip the
destination-disk-too-small check). They do not. \texttt{-k1} operates on
the partition layout but cannot move metadata blocks already laid out at
high offsets in the source ext4 filesystem; \texttt{-icds} only
suppresses the up-front check, it does not prevent the runtime seek
failure.

\textbf{The fix.} The partition layout itself must be made physically
smaller than the smallest target disk, on a copy of the master, before
the image is recaptured. The four-step procedure is mandatory and
order-sensitive:

\begin{enumerate}
\def\labelenumi{\arabic{enumi}.}
\tightlist
\item
  \texttt{e2fsck\ -fy\ /dev/\textless{}source\textgreater{}} --- force a
  clean filesystem before resizing. The resize tools refuse to operate
  on a dirty filesystem.
\item
  \texttt{resize2fs\ /dev/\textless{}source\textgreater{}\ 20G} ---
  shrink the \textbf{filesystem} to a size chosen to be larger than the
  actual data (12 GB → 20 GB gives margin) but smaller than the smallest
  target disk (238 GB).
\item
  \texttt{parted\ /dev/\textless{}sourcedisk\textgreater{}\ resizepart\ \textless{}N\textgreater{}\ 22100MB}
  --- shrink the \textbf{partition} to slightly larger than the
  filesystem, accounting for the partition start offset.
\item
  \texttt{e2fsck\ -fy\ /dev/\textless{}source\textgreater{}} again ---
  verify that the filesystem survived the partition shrink.
\end{enumerate}

Then the image is recaptured with \texttt{partclone} directly
(\texttt{partclone.vfat} for the EFI system partition,
\texttt{partclone.ext4} for the root partition), and the partition-table
dumps are regenerated with \texttt{parted}, \texttt{sgdisk},
\texttt{sfdisk}, and \texttt{blkid}. The result is a smaller image
(\textasciitilde3.6 GB compressed) that restores cleanly to both 238 GB
SSDs and 466 GB HDDs.

The work happens once, off-line, on a workstation in Barcelona, against
a \texttt{qcow2} copy of the master disk attached via \texttt{qemu-nbd};
doing it on-site against the original master is risky and slow. The
recipe in \texttt{3-Guide/Laptop-Deployment/index.md} codifies the
off-line workflow as \textbf{Phase 3} of the deployment and gates it
behind a clearly-marked ``skip if all disks are the same size''
admonition so operators do not incur the work for deployments that do
not need it.

The order of operations is the part that bites. Shrinking the partition
\emph{before} the filesystem truncates the filesystem and destroys data;
the recipe states this in a \texttt{!!!\ warning} box for the same
reason as the matching warning in
\texttt{Wireless-Mesh/1-Static-IP-Mesh/}: a single sentence at the right
place saves a multi-hour recovery later.

\hypertarget{b.6-pxe-server-architecture}{%
\subsubsection{3.B.6 PXE server
architecture}\label{b.6-pxe-server-architecture}}

A PXE deployment runs three services on a single host on the isolated
deployment subnet of §3.A.7:

\begin{longtable}[]{@{}
  >{\raggedright\arraybackslash}p{(\columnwidth - 6\tabcolsep) * \real{0.2500}}
  >{\raggedright\arraybackslash}p{(\columnwidth - 6\tabcolsep) * \real{0.2500}}
  >{\raggedright\arraybackslash}p{(\columnwidth - 6\tabcolsep) * \real{0.2500}}
  >{\raggedright\arraybackslash}p{(\columnwidth - 6\tabcolsep) * \real{0.2500}}@{}}
\toprule\noalign{}
\begin{minipage}[b]{\linewidth}\raggedright
Service
\end{minipage} & \begin{minipage}[b]{\linewidth}\raggedright
Port
\end{minipage} & \begin{minipage}[b]{\linewidth}\raggedright
Package
\end{minipage} & \begin{minipage}[b]{\linewidth}\raggedright
Role
\end{minipage} \\
\midrule\noalign{}
\endhead
\bottomrule\noalign{}
\endlastfoot
DHCP & 67 & \texttt{isc-dhcp-server} & Issue IP, point client at the
boot file \\
TFTP & 69 & \texttt{tftpd-hpa} & Serve GRUB EFI binary, kernel,
initrd \\
NFS & 2049 & \texttt{nfs-kernel-server} & Serve the Clonezilla rootfs
and the image directory \\
\end{longtable}

The PXE host is itself one of the deployed laptops or a small mini-PC;
it does not need server-class hardware. In Gochas, one of the ThinkPads
earmarked for the school was reassigned as PXE host for the duration of
the imaging session and then re-imaged from the same deployment as its
last action.

\textbf{GRUB EFI generation.} The bootloader served via TFTP is
generated by
\texttt{grub-mknetdir\ -\/-net-directory=/tftpboot/nbi\_img\ -\/-subdir=/grub},
which emits \texttt{core.efi} plus the GRUB modules and font. The EFI
binary is then copied as \texttt{bootx64.efi} at the TFTP root because
that is the filename UEFI firmware expects when a DHCP \texttt{filename}
option is presented for a UEFI client (DHCP option 93 architecture
\texttt{00:07} or \texttt{00:09}).

\textbf{Two copies of the kernel, on purpose.} The Clonezilla
\texttt{vmlinuz} and \texttt{initrd.img} exist in \emph{two}
directories: under \texttt{/tftpboot/nbi\_img/} (served via TFTP,
fetched by GRUB during boot) and under \texttt{/tftpboot/clonezilla/}
(served via NFS, mounted by the running kernel as its root). The
duplication is necessary because \texttt{tftpd-hpa} runs in
\texttt{-\/-secure} mode, which chroots the TFTP server to its root
directory and refuses to follow symlinks pointing outside; placing the
files where each service can actually serve them is simpler and more
robust than trying to defeat the chroot.

\textbf{DHCP scope.} A small \texttt{/24} is enough
(\texttt{10.0.0.0/24} in Gochas), with the PXE host on \texttt{.1}, a
short DHCP range (\texttt{.101}--\texttt{.120}) for clients, and a
\texttt{next-server} pointer back to \texttt{.1}. The
architecture-conditional \texttt{filename} directive is the part that
makes the same DHCP scope work for both UEFI and legacy-BIOS clients
without re-configuration.

\hypertarget{b.7-auto-restore-script-and-secure-boot-handling}{%
\subsubsection{3.B.7 Auto-restore script and Secure Boot
handling}\label{b.7-auto-restore-script-and-secure-boot-handling}}

Once GRUB loads, the operator should not be asked to make any choices.
The default GRUB menu entry calls a single \texttt{auto-restore.sh}
script served from the NFS-exported image directory. The script:

\begin{enumerate}
\def\labelenumi{\arabic{enumi}.}
\item
  \textbf{Detects the target disk.} Probes \texttt{/dev/nvme0n1} then
  \texttt{/dev/sda} then \texttt{/dev/vda}, in that order, and selects
  the first present block device. This is what makes the same image and
  the same kernel command line work across machines with NVMe SSDs and
  SATA HDDs without per-machine configuration.
\item
  \textbf{Invokes \texttt{ocs-sr}} with the flags the recipe pins:

\begin{verbatim}
ocs-sr -g auto -e1 auto -e2 -r -j2 -icds -k1 -scr -p reboot \
  restoredisk "$IMAGE_NAME" "$DISK"
\end{verbatim}

  The relevant flags are \texttt{-k1} (proportional partition resize),
  \texttt{-icds} (skip destination-size check --- safe now that the
  image of §3.B.5 fits), \texttt{-scr} (skip restorability check), and
  \texttt{-p\ reboot} (reboot when done so the operator only has to
  power-cycle once).
\item
  \textbf{Reboots into the deployed OS.} The newly-imaged disk is now
  the first boot device and Linux Mint comes up with the \texttt{aucoop}
  user ready to log in.
\end{enumerate}

\textbf{Secure Boot.} The single biggest field-debugging cost across the
deployments was the silent failure mode of UEFI firmware presented with
an unsigned \texttt{bootx64.efi}: the firmware downloads the binary,
\emph{silently} discards it, and falls through to the next entry in the
boot order (typically IPv6 PXE, which times out). There is no error
message on screen and the only diagnostic clue is the silence itself.
The recipe gates the entire deployment on a one-time BIOS step (Step 14:
disable Secure Boot on every target machine before imaging) and the
lessons list of §3.B.9 elevates this to an inventory-time check in
§3.B.2 so the failure cannot recur. Secure Boot may be re-enabled after
deployment if the receiving institution's policy demands it; the
deployed Linux Mint does support it.

\hypertarget{b.8-quality-control-and-user-handover}{%
\subsubsection{3.B.8 Quality control and user
handover}\label{b.8-quality-control-and-user-handover}}

The last step before a machine leaves the imaging table is a short
quality check, performed in the order in which the failures it catches
typically appear:

\begin{enumerate}
\def\labelenumi{\arabic{enumi}.}
\tightlist
\item
  The machine boots from the local disk to the Linux Mint login screen
  unaided (no PXE, no USB).
\item
  The \texttt{aucoop} user logs in with the documented password.
\item
  Wi-Fi can associate to the production network of §3.A.5 and reach the
  internet.
\item
  OnlyOffice opens a sample document and renders it correctly.
\item
  The hostname matches the inventory record of §3.B.2.
\end{enumerate}

A failed check sends the machine back to the corresponding step (a boot
failure to §3.B.5, an image-content failure to §3.B.3, a network failure
to §3.A.5) rather than triggering an ad-hoc fix on the individual
machine. The discipline is what keeps a fleet of nine or twelve laptops
actually identical at handover.

Handover to the receiving institution adds a short briefing --- how to
log in, how to reach Wi-Fi, where the printed runbook is, who to contact
--- and a written record that the machine has been delivered. End-user
training at depth is out of scope for this thesis; what is in scope is
delivering machines whose default configuration does not require that
training to be useful.

\hypertarget{b.9-lessons-learned-and-inputs-to-the-handbook}{%
\subsubsection{3.B.9 Lessons learned and inputs to the
handbook}\label{b.9-lessons-learned-and-inputs-to-the-handbook}}

As in §3.A, the laptop deployments produced a set of lessons that landed
directly in the handbook. They are listed here because each one points
to a specific change the next operator should not have to re-discover.

\textbf{Secure Boot is an inventory-time check, not a debug-time
discovery.} The hours lost to silent UEFI rejection of an unsigned GRUB
binary turned this from a recipe footnote into the gating BIOS posture
of §3.B.2 and a \texttt{!!!\ warning} admonition at the top of the PXE
recipe.

\textbf{The partition-resize problem deserves its own phase.} The
\texttt{target\ seek\ ERROR} failure was the reason
\texttt{3-Guide/Laptop-Deployment/} was reorganised into four explicit
phases (prepare, capture, resize, deploy) rather than a single linear
list of steps. Operators with uniform-disk fleets skip Phase 3 entirely;
operators with mixed disks know exactly which phase to read.

\textbf{\texttt{tftpd-hpa\ -\/-secure} does not follow symlinks.} This
single sentence is now a \texttt{!!!\ warning} in the recipe; without
it, an operator who tries to keep one canonical copy of the Clonezilla
files and symlink the rest sees TFTP serve nothing and has no obvious
diagnostic.

\textbf{Auto-detect the target disk.} Hard-coding \texttt{/dev/sda} in
the restore command works on every ThinkPad with a SATA disk and fails
on every ThinkPad with an NVMe SSD. The auto-detection script in §3.B.7
was written after the first NVMe machine refused to image; it is now the
default in the recipe.

\textbf{Use \texttt{-k1\ -icds\ -scr\ -p\ reboot} together with a
resized image, not instead of one.} These flags relax \texttt{ocs-sr}'s
safety checks but do not fix the underlying ext4 layout problem of
§3.B.5. The recipe states this explicitly so the next operator does not
lose a day chasing flag combinations that cannot work.

\textbf{A simple stack beats DRBL for this fleet size.} A full DRBL
deployment is overkill for nine or twelve machines and adds a layer of
debugging that an on-site team cannot afford. The recipe is built on
plain \texttt{isc-dhcp-server} + \texttt{tftpd-hpa} +
\texttt{nfs-kernel-server} + Clonezilla Live precisely so that every
component can be inspected and restarted independently when something
misbehaves. DRBL becomes attractive at a different fleet size; that is a
future-work note in §7.

\textbf{The image is a versioned artefact.} Naming the image directory
\texttt{aucoop-mint22.3-2026-03} and storing both the \texttt{qcow2}
master copy and the resized image in the project repository is what
allows §4 to claim that \emph{this specific image} was deployed at
\emph{this specific site on this specific date}. The companion thesis
{[}Motje, 2026{]} inherits the same naming convention for service-side
artefacts.

These lessons feed into the §4 validation against the three lenses
(coverage: did every target laptop receive the image?; sufficiency: did
the resulting fleet meet the school's actual usage in the days that
followed?; adaptation: could a second team in a different site execute
the same recipe?) and into the per-objective check of §7.

\begin{center}\rule{0.5\linewidth}{0.5pt}\end{center}

\hypertarget{c-the-aucoop-handbook-as-a-knowledge-artefact}{%
\subsection{3.C The AUCOOP Handbook as a knowledge
artefact}\label{c-the-aucoop-handbook-as-a-knowledge-artefact}}

The two preceding parts of this chapter (§3.A and §3.B) describe
\emph{what} was deployed in the field. This third part describes
\emph{how the deployment knowledge itself is preserved} so that the next
student, the next volunteer, or the next partner association does not
have to start again from a blank page. It is the most original
contribution of this thesis: not a single network or a single classroom
of refurbished laptops, but a reusable instrument --- the
\textbf{Community Network Handbook} --- capable of producing many of
them.

The handbook is hosted at
\url{https://github.com/aucoop/Community-Network-Handbook} and published
as a static website plus a downloadable PDF. The thesis writes
\emph{about} the handbook; the handbook itself is the durable
deliverable.

\hypertarget{c.1-why-a-living-handbook}{%
\subsubsection{3.C.1 Why a living
handbook}\label{c.1-why-a-living-handbook}}

§1.2 framed the problem of \emph{knowledge volatility} in a volunteer
association. AUCOOP runs on bachelor and master students who join,
contribute for one or two academic years, graduate, and leave. Every
project --- Namibia 2024, Namibia 2026, the local pilots in Barcelona
--- produced internal documents (PDFs, slide decks, hand-written field
notes) that ended up in shared drives that nobody opens once the
contributors are gone. The next team rediscovers the same DHCP option,
the same OpenWrt mesh quirk, the same Clonezilla
\texttt{partclone\ target\ seek\ ERROR}, and pays the same debugging
tax.

A living handbook reverses this dynamic by making three structural
commitments:

\begin{enumerate}
\def\labelenumi{\arabic{enumi}.}
\tightlist
\item
  \textbf{The artefact is the documentation, not a side-effect of it.}
  Contributors write directly into the handbook during the project, not
  into private notes that \emph{might} be transcribed afterwards.
\item
  \textbf{The artefact is publicly readable.} External collaborators
  (NGO partners, future host schools, other student associations) can
  consult it without asking permission, which removes the friction that
  kills internal wikis.
\item
  \textbf{The artefact is editable through a low-ceremony,
  version-controlled workflow} (Git + Markdown + pull request), so
  corrections from the field are cheap to land.
\end{enumerate}

Compared with the realistic alternatives --- a Google Drive folder, a
Notion workspace, a wiki on a self-hosted server that nobody patches ---
the chosen model trades convenience (no rich WYSIWYG editor) for
longevity (plain text under version control survives platform changes
and credential losses).

\hypertarget{c.2-information-architecture}{%
\subsubsection{3.C.2 Information
architecture}\label{c.2-information-architecture}}

The handbook is organised in four top-level chapters, each with a clear
epistemic role:

\begin{longtable}[]{@{}
  >{\raggedright\arraybackslash}p{(\columnwidth - 4\tabcolsep) * \real{0.3333}}
  >{\raggedright\arraybackslash}p{(\columnwidth - 4\tabcolsep) * \real{0.3333}}
  >{\raggedright\arraybackslash}p{(\columnwidth - 4\tabcolsep) * \real{0.3333}}@{}}
\toprule\noalign{}
\begin{minipage}[b]{\linewidth}\raggedright
Chapter
\end{minipage} & \begin{minipage}[b]{\linewidth}\raggedright
Role
\end{minipage} & \begin{minipage}[b]{\linewidth}\raggedright
Voice
\end{minipage} \\
\midrule\noalign{}
\endhead
\bottomrule\noalign{}
\endlastfoot
\texttt{1-Introduction} & Why the handbook exists, who it is for &
Editorial \\
\texttt{2-Imaginary-Use-Case} & A fictional community network, told as a
story; one section per challenge & Narrative, second-person, italic
question titles \\
\texttt{3-Guide} & Step-by-step recipes; one folder per technology &
Instructional, imperative \\
\texttt{4-Real-Use-Cases} & Concrete deployments (Namibia, \ldots) used
as case studies & Reporting \\
\end{longtable}

The defining design rule is the \textbf{1-to-1 mapping between Chapter 2
and Chapter 3}: every story section in Chapter 2 must have a matching
recipe in Chapter 3, and vice versa, with explicit cross-links in both
directions. The intent is twofold. A reader who is \emph{learning the
domain} enters through the story and follows the link to the recipe when
they want to act. A reader who is \emph{executing a deployment} enters
through the recipe and follows the link to the story when they need to
understand the trade-offs. Neither audience is penalised. The constraint
is enforced at review time and is mechanical enough that it can be
audited (see §3.C.5).

Inside a section folder, content is co-located with its images:

\begin{verbatim}
docs/3-Guide/Laptop-Deployment/
├── index.md
├── AUCOOP-image.md
└── images/
    ├── pxe-architecture.webp
    └── clonezilla-savedisk.webp
\end{verbatim}

Image filenames are prefixed with the owning section
(\texttt{Cudy-WR3000E-luci.webp}, \texttt{2.1-router-front-panel.webp})
so that a file moved or renamed without its section is immediately
recognisable as orphaned. All raster images use WebP at quality 85,
chosen as the smallest format that still renders well in the PDF build.

\hypertarget{c.3-toolchain}{%
\subsubsection{3.C.3 Toolchain}\label{c.3-toolchain}}

The handbook is built with \textbf{Zensical}, a static-site generator
that consumes Markdown plus a \texttt{mkdocs.yml} navigation manifest
and produces both the website and the PDF. The relevant configuration
choices are:

\begin{itemize}
\tightlist
\item
  \textbf{Theme.} Material for MkDocs in \texttt{navigation.sections}
  mode, with light/dark/system palette toggles. The teal accent matches
  the AUCOOP branding.
\item
  \textbf{Markdown extensions.} \texttt{pymdownx.superfences} (with a
  Mermaid custom fence), \texttt{admonition}, \texttt{pymdownx.details},
  \texttt{pymdownx.highlight}, \texttt{pymdownx.tabbed},
  \texttt{pymdownx.arithmatex}, \texttt{attr\_list}, \texttt{tables},
  \texttt{md\_in\_html}. These cover the four expressive needs of a
  deployment guide: diagrams, side notes, code with line numbers, and
  tabular data.
\item
  \textbf{Diagram pipeline.} Mermaid is rendered client-side from fenced
  code blocks, which keeps diagrams in plain text inside the source
  repository and reviewable in pull requests. PlantUML is available for
  the rarer cases that need it, served by \texttt{build\_plantuml}
  against the public PlantUML server.
\item
  \textbf{PDF target.} The repository's release pipeline produces a
  single \texttt{Community-Network-Handbook.pdf} artefact; the website
  surfaces it through the \texttt{extra.book\_download\_url} setting.
  Both outputs are generated from the same Markdown source, so there is
  no risk of the printable copy and the web copy drifting apart.
\end{itemize}

The build commands are deliberately wrapped: contributors run
\texttt{zensical\ serve} for local preview and \texttt{zensical\ build}
for production, never \texttt{mkdocs} directly. Documenting this in the
project's \texttt{AGENTS.md} removes a common source of confusion when a
new contributor's local environment behaves differently from CI.

\hypertarget{c.4-dual-output-web-and-pdf}{%
\subsubsection{3.C.4 Dual output --- web and
PDF}\label{c.4-dual-output-web-and-pdf}}

A handbook intended for low-connectivity contexts cannot be a
website-only artefact. During the Gochas deployment, the school's uplink
was unusable for hours at a time, and the team relied on the offline PDF
copy carried on a laptop. Conversely, an NGO evaluating AUCOOP from
Barcelona benefits from a searchable, linkable web version. The dual
output is not a redundancy; it serves two distinct usage modes:

\begin{itemize}
\tightlist
\item
  The \textbf{web} version is the canonical, evolving document. It is
  search-indexed, link-friendly, and updated on every merge to
  \texttt{main}.
\item
  The \textbf{PDF} version is the \emph{expedition copy}. It is produced
  as a GitHub release artefact, pinned to a specific commit, and
  intended to be downloaded before travelling.
\end{itemize}

Pinning the PDF to a release tag means a field team can reference
\emph{the exact version of the handbook they took with them} months
later, even if the website has moved on. This addresses a classic
reproducibility problem in deployment documentation: the recipe that
worked in March may have been silently rewritten by July.

\hypertarget{c.5-contribution-model-and-ai-assisted-authoring}{%
\subsubsection{3.C.5 Contribution model and AI-assisted
authoring}\label{c.5-contribution-model-and-ai-assisted-authoring}}

The handbook is written by a rotating set of contributors of unequal
experience: senior PhD students, master students writing their thesis
(this thesis among them), and bachelor students contributing for a
single semester. The contribution model has to be strict enough that
quality does not collapse between cohorts, but light enough that a
one-semester contributor can land a useful change in their first week.
Three mechanisms make this work.

\textbf{Rule files.} The conventions that govern the handbook are not
folklore; they are version-controlled Markdown files under
\texttt{.opencode/rules/}. Four files cover, respectively, the general
project rules (\texttt{general.md}), the narrative voice required for
Chapter 2 stories (\texttt{chapter2-story.md}), the imperative structure
required for Chapter 3 recipes (\texttt{chapter3-guide.md}), and the
navigation-manifest invariants for \texttt{mkdocs.yml}
(\texttt{mkdocs-nav.md}). A new contributor is told to read these once;
a reviewer who finds a violation points to the exact rule. The rules
also encode operational details that are easy to forget --- the
placeholder image workflow, the WIP admonition format, the file-naming
conventions for sections and images.

\textbf{OpenCode subagents.} The repository ships with five project
subagents defined in \texttt{.opencode/agents/} that can be invoked from
any OpenCode-compatible editor:

\begin{longtable}[]{@{}
  >{\raggedright\arraybackslash}p{(\columnwidth - 4\tabcolsep) * \real{0.3333}}
  >{\raggedright\arraybackslash}p{(\columnwidth - 4\tabcolsep) * \real{0.3333}}
  >{\raggedright\arraybackslash}p{(\columnwidth - 4\tabcolsep) * \real{0.3333}}@{}}
\toprule\noalign{}
\begin{minipage}[b]{\linewidth}\raggedright
Agent
\end{minipage} & \begin{minipage}[b]{\linewidth}\raggedright
Role
\end{minipage} & \begin{minipage}[b]{\linewidth}\raggedright
Mode
\end{minipage} \\
\midrule\noalign{}
\endhead
\bottomrule\noalign{}
\endlastfoot
\texttt{@writer} & Create or expand sections; runs against the rule
files; can edit and shell out & Read--write \\
\texttt{@reviewer} & Quality review of an existing section against the
rules & Read-only \\
\texttt{@structure} & Refactor folder structures and synchronise
\texttt{mkdocs.yml} nav & Read--write \\
\texttt{@diagrams} & Author or improve Mermaid diagrams & Read--write \\
\texttt{@consistency} & Audit cross-cutting invariants (Ch2↔Ch3 mapping,
broken links, missing TODOs) & Read-only \\
\end{longtable}

The \texttt{@writer} agent is the most-used. Its rule of operation is
informative: it must read the relevant rule file \emph{before} editing,
present a compact plan, wait for explicit human approval (unless invoked
through a custom command, which counts as approval), and only then
write. This pattern --- \emph{plan, approve, execute} --- keeps the
human reviewer in the loop on every change while absorbing the
mechanical cost of conforming to the rules.

\textbf{Custom commands.} Five slash-commands package recurring
workflows so that they are issued the same way every time:

\begin{itemize}
\tightlist
\item
  \texttt{/new-section} --- create a paired Ch2 story stub and Ch3
  recipe stub, plus the \texttt{mkdocs.yml} nav entries.
\item
  \texttt{/review-chapter} --- invoke \texttt{@reviewer} against a
  chapter and produce a review report.
\item
  \texttt{/add-diagram} --- drop a Mermaid block in the right place with
  the right fencing.
\item
  \texttt{/audit} --- run \texttt{@consistency} across the whole
  repository.
\item
  \texttt{/guide-from-steps} --- turn a contributor's raw shell history
  or field notes into a Chapter 3 recipe that conforms to the structural
  template.
\end{itemize}

The combination is what enables a master-thesis contribution like this
one to add four merged feature branches over a few weeks while keeping
the handbook's voice and structure intact. The cost of consistency is
paid by tooling, not by reviewer fatigue.

It should be stated openly that AI assistance is part of the workflow.
The handbook is not generated by AI --- every page is reviewed and
approved by a human contributor --- but the mechanical conformance work
(filling templates, synchronising the nav, regenerating tables, applying
tone rules) is delegated to subagents. This is itself a contribution: it
documents a viable model for small-association open documentation
projects that have neither paid editors nor a stable maintainer pool.

\hypertarget{c.6-governance-for-continuity}{%
\subsubsection{3.C.6 Governance for
continuity}\label{c.6-governance-for-continuity}}

Tooling alone does not guarantee survival. The handbook also needs an
explicit governance contract. The following elements are proposed and
are in force on the \texttt{dev\_mj\_thesis} branch at the time of
writing:

\begin{itemize}
\tightlist
\item
  \textbf{Maintainer rotation.} At any time the handbook has at least
  one \emph{maintainer} (currently the AUCOOP project lead) and at least
  one \emph{active contributor} whose thesis depends on the project. The
  active contributor is expected to take over maintainer duties at the
  end of their thesis if no successor has joined. A short hand-over
  document accompanies each transition.
\item
  \textbf{Definition of Done for a chapter.} A chapter is considered
  ``done for this iteration'' when (i) every section has at least one
  paragraph of body text, (ii) the Ch2↔Ch3 cross-links exist in both
  directions, (iii) all images are real (no placeholders), (iv)
  \texttt{/audit} reports zero structural errors, and (v)
  \texttt{zensical\ build} succeeds without warnings.
\item
  \textbf{Picking up an open chapter.} A new contributor opens the
  chapter folder, reads its \texttt{index.md}, runs \texttt{/audit} to
  see the outstanding TODOs and WIP markers, and chooses one. The
  combination of \texttt{!!!\ info\ "Work\ in\ Progress"} admonitions
  and \texttt{\textless{}!-\/-\ TODO:\ ...\ -\/-\textgreater{}} comments
  is dual on purpose: the first is visible to readers of the rendered
  site and signals that they should not rely on the section yet; the
  second is invisible to readers but \texttt{grep}-able by contributors.
\end{itemize}

This governance is light by design. Heavier processes --- RFC documents,
formal release schedules, CODEOWNERS --- would not survive the next
volunteer rotation. The chosen mechanisms degrade gracefully: even if
the rules and agents fall out of use for a semester, the resulting
damage is limited to inconsistent voice, not a broken build or a
corrupted history.

\hypertarget{c.7-what-this-contributes-beyond-the-deployment}{%
\subsubsection{3.C.7 What this contributes beyond the
deployment}\label{c.7-what-this-contributes-beyond-the-deployment}}

The Namibia deployment described in §4 will eventually be one row in the
handbook's \texttt{4-Real-Use-Cases} chapter. The handbook itself will
outlive it. The contribution claimed here is therefore not ``documenting
Namibia'' but \emph{producing the instrument that makes Namibia
documentable in a form the next team can act on}. The thesis treats this
instrument with the same rigour as the network and endpoint work: as a
designed system with explicit requirements (continuity, dual output, low
contributor onboarding cost), explicit components (rule files,
subagents, custom commands, governance contract), and a measurable
acceptance criterion (the four merged branches of
\texttt{dev\_mj\_thesis}, which collectively land a complete first
version of the hardware-network and laptop-deployment material).
